% Paper D, Journal of Automated Reasoning version. Edited directly since the JAR revision of
% 2026-11-01 (branch jar-revision); port.py and FIDELITY.md describe the earlier generated
% version and are no longer used (paper/README.md).
%
% Template: Springer Nature LaTeX authoring template (sn-jnl.cls, December 2024 package),
% reference style sn-mathphys-num (Math and Physical Sciences, numbered).
%
% Revision tags: a comment line "% [rev X]" precedes every paragraph that the JAR revision
% added or changed, X being the item of paper/jar/CHANGES-FOR-APPROVAL.md (D1-D7, m1-m13, SD
% for the Statements and Declarations). They can be deleted once the author has approved the
% changes.
%
% Open before submission: every [AUTHOR: ...] question and every [AUTHOR TO APPROVE]
% placeholder (listed in CHANGES-FOR-APPROVAL.md).

\documentclass[pdflatex,sn-mathphys-num]{sn-jnl}

\usepackage{amsmath,amssymb,amsfonts}
\usepackage{amsthm}
\usepackage{listings}
\usepackage{array}
\usepackage[capitalise,noabbrev]{cleveref}
\newcolumntype{L}[1]{>{\raggedright\arraybackslash}p{#1}}
\newcolumntype{R}[1]{>{\raggedleft\arraybackslash}p{#1}}

% Theorem environments: one shared counter (Definition 1, Theorem 2).
\theoremstyle{thmstyleone}
\newtheorem{theorem}{Theorem}
\theoremstyle{thmstylethree}
\newtheorem{definition}[theorem]{Definition}
\crefname{definition}{Definition}{Definitions}
\crefname{lstlisting}{Listing}{Listings}
\newcommand{\crefrangeconjunction}{--}% "Sections 4--7"

% [rev D5] The mark of a claim argued on paper only.
\newcommand{\informal}{\ensuremath{^{\dagger}}}

% sn-jnl's \subparagraph is a run-in heading, which a listing that follows at once would
% swallow; print the pending heading before the listing (layout only).
\makeatletter
\AddToHook{env/lstlisting/before}{\if@noskipsec\leavevmode\fi}
\makeatother

% Appendices are referred to as "Appendix A".
\AtBeginEnvironment{appendices}{\crefalias{section}{appendix}}

\lstdefinelanguage{lean}{
  morekeywords={theorem,def,abbrev,lemma,structure,where,fun,let,if,then,else,instance,
    noncomputable,by,exact,refine,unfold,have,obtain},
  sensitive=true,
  morecomment=[l]{--},
  morestring=[b]",
}
\lstset{
  language=lean,
  basicstyle=\ttfamily\small,
  keywordstyle=\bfseries,
  columns=fullflexible,
  keepspaces=true,
  xleftmargin=1em,
  breaklines=true,
  breakindent=2em,
  literate=
    {ℕ}{{$\mathbb{N}$}}1
    {∀}{{$\forall$}}1
    {×}{{$\times$}}1
    {∃}{{$\exists$}}1
    {∈}{{$\in$}}1
    {∉}{{$\notin$}}1
    {∧}{{$\wedge$}}1
    {∨}{{$\vee$}}1
    {↔}{{$\leftrightarrow$}}1
    {→}{{$\to$}}1
    {↪}{{$\hookrightarrow$}}1
    {≤}{{$\le$}}1
    {<}{{$<$}}1
    {≠}{{$\neq$}}1
    {∅}{{$\emptyset$}}1
    {ᶜ}{{$^{c}$}}1
    {·}{{$\cdot$}}1
    {⟨}{{$\langle$}}1
    {⟩}{{$\rangle$}}1
    {⦃}{{$\{\!|$}}1
    {⦄}{{$|\!\}$}}1
    {₀}{{$_0$}}1
    {ᵢ}{{$_i$}}1
    {ⱼ}{{$_j$}}1
    {φ}{{$\varphi$}}1
    {α}{{$\alpha$}}1
    {β}{{$\beta$}}1
    {ι}{{$\iota$}}1
    {𝒱}{{$\mathcal{V}$}}1
    {𝒟}{{$\mathcal{D}$}}1
    {∖}{{$\setminus$}}1
    {∩}{{$\cap$}}1
    {…}{{$\dots$}}1
    {₁}{{$_1$}}1
    {₂}{{$_2$}}1
}


\raggedbottom

\begin{document}

\title[Machine-Checked Complexity of Decision-Table Completeness]{Machine-Checked Complexity of Decision-Table Completeness in Lean}

\author*[1]{\fnm{Carlos A.} \sur{Russell}}\email{carlos.a.russell@gmail.com}
% ORCID: https://orcid.org/0009-0003-3741-0230
% (sn-jnl prints no ORCID; Springer collects it in the submission system.)

\affil*[1]{\orgname{Independent Researcher}, \orgaddress{\city{Buenos Aires}, \country{Argentina}}}

% [rev D6]
\abstract{%
A decision table is complete when every combination of its declared condition values that the
domain constraints allow is covered by a row. We prove
in Lean~4, over multi-tape Turing machines and with no unproved assumptions beyond the
definitions themselves, that for finite-domain decision tables with clause constraints and rows
that are complete condition tuples, deciding whether some combination is missing is
NP-complete, deciding completeness is coNP-complete, and counting the missing combinations is
\#P-complete, all already for Boolean domains and empty tables. An adequacy theorem ties these
word problems to the set-based definitions of an existing formalisation of decision-table
completion, so the results speak about the original objects. Complementary polynomial-time
results locate the hardness: all three problems are polynomial without constraints, and with a
fixed number of attributes, where one machine runs within $c\,(|x|+1)^{en+e}$ on instances with
$n$ attributes, for some constants $c$ and~$e$. For monotone two-literal constraints, counting
the missing combinations remains \#P-complete under reductions with a single oracle call; for
this we give a new formal proof, in the same Turing-machine model, of a one-call version of
Valiant's theorem on \#Monotone-2SAT. We also report how the development was engineered and how
an AI system was used to write it. The classification was first argued on paper in the
author's earlier work, which is under review; every proof here is self-contained.}

% [rev D6] Six keywords (JAR asks for 4-6).
\keywords{decision tables, completeness, coNP-completeness, \#P-completeness, Lean 4,
  Cook--Levin theorem}

% ACM 2012 classification, kept for the submission form:
%   Theory of computation~Problems, reductions and completeness (500);
%   Theory of computation~Logic and verification (300).

\maketitle

%-----------------------------------------------------------------------------------------------
\section{Introduction}
\label{sec:intro}

Decision tables state business rules, regulations and clinical guidelines as rows of
condition values~\cite{PollackHicksHarrison1971}. A table is \emph{complete} when every combination of declared condition
values that the domain constraints allow is covered by some row; the uncovered combinations
form its \emph{missing set}~$M(T_0)$. Completeness checking is a staple of decision-table
practice, from early expert-system checkers~\cite{SuwaScottShortliffe1982} to the analysers of
the DMN standard~\cite{CalvaneseDLMMT2016}, which report missing rules. Counting them matters
too: one missing edge case and a table that misses half of its valid inputs are both
``incomplete'', and $|M(T_0)|$ is what tells them apart.

% [rev D6]
Russell~\cite{PaperA}, in a paper under review, re-examined a decision-table completion method
over declared finite domains and argued, on paper, that completeness under clause constraints
is coNP-complete and that counting $|M(T_0)|$ is \#P-complete (Proposition~5 there). The
present paper does not depend on that argument: every proof it gives is its own. The
accompanying Lean development, latc-check~\cite{latc_check_0_1_0}, formalised the
decision-table model and a combinatorial reduction from its own propositional formulas. It did
not formalise the complexity classes, which need a model of computation, polynomial-time bounds
and encodings.

This paper supplies them. We compose latc-check's model with Complexitylib~\cite{Complexitylib},
a Lean~4 library of complexity theory on multi-tape Turing machines with a machine-checked
Cook--Levin theorem, and prove the classification end to end, together with results that neither
earlier development contains. On paper the classification is elementary: with no rows,
completeness is unsatisfiability of the constraints. The contribution is its mechanisation on
Turing machines with an adequacy theorem, the tractability boundary with an explicit degree, and
the one-call Valiant theorem. Every theorem we cite depends only on Lean's three standard
axioms, and the build contains no \lstinline{sorry}.

\emph{Scope.} Throughout, a row of the table is a complete tuple of condition values, as in
latc-check. Tables with compact rows, such as DMN rules with wildcards or ranges, are a different
problem: there completeness is coNP-hard even without constraints (\Cref{sec:related}).

% [rev D5]
\emph{Informal claims.} A few claims are argued on paper only and are not theorems of the
development. They are marked with a dagger, as in ``P\informal'', both in the text and in
\Cref{tab:boundary}. Every other result is a Lean theorem, cited by its name.

\subparagraph*{Contributions.}
\Cref{tab:contrib} separates what this paper adds from what it reuses.
\begin{enumerate}
% [rev D5]
\item \textbf{The classification, machine-checked end to end} (\Cref{sec:decision,sec:counting}):
  non-emptiness of $M(T_0)$ is NP-complete, completeness coNP-complete and counting \#P-complete,
  for finite domains with clause constraints, with every number of the input written in binary.
  Hardness holds already for instances with Boolean domains and empty tables, a Lean theorem
  about the shape of the reduction (\Cref{sec:decision}).
\item \textbf{Adequacy} (\Cref{sec:encoding}): the encoded problems are exactly latc-check's
  set-based ones.
% [rev D7]
\item \textbf{One-call hardness for monotone two-literal counting} (\Cref{sec:monotone}): a
  definition of one-call reductions in Complexitylib's model, and a new formal proof in that
  model of a one-call version of Valiant's theorem, by classical techniques. From it, counting
  the missing combinations stays \#P-complete under one-call reductions when every constraint is
  $A_i = 1 \lor A_j = 1$, for tables with arbitrary rows (hard already with none).
\item \textbf{The tractability boundary} (\Cref{sec:boundary}): polynomial time without
  constraints, and with a fixed number of attributes, for both decision problems and for
  counting, with an explicit degree linear in the number of attributes.
% [rev D1]
\item \textbf{An account of the engineering} (\Cref{sec:engineering}): the size of each part,
  what Complexitylib provided and lacked, how the polynomial-time proofs were automated, how the
  instance-dependent bounds were kept explicit, the lemmas that could be reused, and how an AI
  system was used.
\end{enumerate}

% [rev D7] Last row of the table.
\begin{table}[t]
\caption{What is reused and what is new.}
\label{tab:contrib}
\centering\small
\begin{tabular}{L{0.38\textwidth}L{0.26\textwidth}L{0.26\textwidth}}
\hline
result & before this paper & here \\
\hline
decision-table model, missing set & latc-check (Lean) & reused \\
SAT reduction to missing states & latc-check, from its own formulas & new, from Complexitylib's CNF \\
NP/coNP-completeness, \#P-completeness & on paper~\cite[Prop.~5]{PaperA} & Turing machines, in Lean \\
adequacy of the encoded problems & -- & new \\
polynomial cases: no constraints, fixed $k$ & analogues for firewalls~\cite{HamiltonRoughanNguyen2022} and for unknown entries~\cite{VuDemetrovics2026} & new for decision tables \\
explicit degree in the number of attributes & -- & new \\
one-call reductions, one-call Valiant & Valiant (Turing reductions); Senellart--Gnatenko (one call, logical model, Lean library) & new formal proof in Complexitylib's model \\
\hline
\end{tabular}
\end{table}

% [rev D1]
\Cref{sec:prelim} recalls Complexitylib and latc-check, \Cref{sec:encoding} fixes the input
language, \Cref{sec:decision,sec:counting,sec:monotone,sec:boundary} prove the results,
\Cref{sec:related} discusses related work, \Cref{sec:formal} the trusted base and
\Cref{sec:engineering} the engineering of the development.
\Cref{app:statements} lists the main Lean statements, \Cref{app:definitions} the trusted
definitions and \Cref{app:artifact} the artifact.

%-----------------------------------------------------------------------------------------------
\section{Preliminaries}
\label{sec:prelim}

\subparagraph*{Machines, encodings and classes.}
Complexitylib, built on CSLib~\cite{CSLib2026} and Mathlib, works with multi-tape Turing
machines over a binary alphabet. A type $\alpha$ is
given an encoding, an injection $\alpha \hookrightarrow \{0,1\}^*$; pairs, lists and natural
numbers have standard encodings, numbers in binary (\lstinline{encNat}) or in unary
(\lstinline{encUnary}). A function $f$ is computable in polynomial time between two encodings
when some machine maps the encoding of every $a$ to the encoding of $f(a)$ within a polynomial
in the input length. A language is in P when its characteristic function is; in NP when it is
the set of inputs having a certificate of polynomially bounded length accepted by a
polynomial-time verifier; and coNP is the class of complements. A function $f$ into
$\mathbb{N}$ is in \#P~\cite{Valiant1979permanent} when $f(a)$ is the number of certificates of length exactly $p(|a|)$,
for a polynomial~$p$, that a polynomial-time verifier accepts on~$a$. A library lemma,
\lstinline{mem_SharpP_of_polyTime_length}, lets the certificate length be any function of the input that is computable in polynomial
time with its value written in unary; our counting verifiers use it.

\subparagraph*{Reductions.}
Three notions are used. A \emph{many-one} reduction from $A$ to $B$ is a polynomial-time $r$
with $a \in A \iff r(a) \in B$; NP-, coNP-hardness and completeness are defined with it. A
\emph{parsimonious} reduction from $f$ to $g$ is a polynomial-time $r$ with $f(a) = g(r(a))$;
Complexitylib defines \#P-hardness with it. These two are Complexitylib's. The third is defined
in this development; it is the polynomial-time \emph{metric reduction} that Krentel defines
between functions into $\mathbb{N}$ for optimisation problems~\cite{Krentel1988}:
\begin{definition}[one-call reduction]
\label{def:onecall}
$f$ reduces to $g$ with one call if there are polynomial-time functions $r$ and
$\mathit{post}$ with $f(a) = \mathit{post}(a, g(r(a)))$ for every~$a$, where $\mathit{post}$
reads $a$ and the answer $g(r(a))$ written in binary. $g$ is \emph{one-call \#P-hard} if every
function in \#P, under any encoding, reduces to it with one call, and \emph{one-call
\#P-complete} if it is moreover in \#P.
\end{definition}
% [rev D5] [rev m12]
The running time of $\mathit{post}$ is polynomial in the length of $a$ and of the oracle's
answer. For a target $g$ whose values have polynomially many bits, as for every function in
\#P, a one-call reduction is a polynomial-time Turing reduction making exactly one oracle query,
whose answer is post-processed together with the input, and one-call hardness implies hardness
under polynomial-time Turing reductions\informal, the notion of
Valiant~\cite{Valiant1979enumeration}. (Without that bound the notion would be too weak: a $g$
with doubly exponential values gives $\mathit{post}$ exponential time in $|a|$ while carrying no
information. We use one-call hardness only for \#Monotone-2SAT, whose value is at most
$2^{\#\text{variables}}$, and for a target proved to be in \#P.) Parsimonious reductions are
one-call, with $\mathit{post}(a, y) = y$ (\lstinline{ParsimoniousReducible.oneCallReducible});
so \#P-hardness implies one-call \#P-hardness (\lstinline{SharpPHard.oneCallSharpPHard}), and a
parsimonious reduction followed by a one-call reduction is one-call
(\lstinline{ParsimoniousReducible.trans_oneCall}).
Zank\'o defines \#P-completeness with a variant of many-one reductions and shows the
$(0,1)$-permanent complete under it~\cite{Zanko1991}. One-call hardness is weaker than parsimonious
hardness, which matters below: a problem whose decision version
is trivial cannot be parsimoniously \#P-hard unless P~=~NP, and can still be one-call
\#P-hard.
Our hardness proofs start from Complexitylib's Cook--Levin theorem and its parsimonious form:
SAT is NP-hard, UNSAT coNP-complete, and \#SAT \#P-hard under parsimonious reductions.

\subparagraph*{latc-check's model.}
% [rev m1] (tuples are written \tau; x is kept for input words)
Attributes are indexed by a finite type~$\iota$, attribute $i$ having the declared finite domain
$D_i$, so the space of condition tuples $\mathcal{D} = \prod_i D_i$ is finite. Constraints are a
decidable predicate $C$ on tuples, and a table is represented by its set $T_0$ of condition
tuples. The valid space is $\mathcal{V} = \{\tau \in \mathcal{D} : C(\tau)\}$, the missing set is
$M(T_0) = \mathcal{V} \setminus T_0$, and the table is complete when $M(T_0) = \emptyset$. In
Lean these are \lstinline{validSpace C} and \lstinline{missing C T}, both \lstinline{Finset}s, so
duplicate rows play no role. The present development reuses this model; it does not use
latc-check's own reduction. Two lemmas check that latc-check's clause constraints and
satisfiability agree with Complexitylib's CNF formulas (\lstinline{holdsAll_eq_eval},
\lstinline{satisfiable_decode_iff}); the hardness proofs below use a new reduction stated
directly on Complexitylib's CNF.

%-----------------------------------------------------------------------------------------------
\section{Instances, encoding and adequacy}
\label{sec:encoding}

\subparagraph*{Instances.}
An instance is a triple $I = (d, S, T)$ of lists of natural numbers:
\begin{itemize}
\item $d = [d_0, \dots, d_{n-1}]$, the domain sizes; attribute $i$ ranges over
  $\{0, \dots, d_i - 1\}$, and $n$, the number of attributes, is the length of~$d$;
\item $S$, the clause constraints: a list of clauses, each a list of conditions $(i, v)$, read
  $A_i = v$;
\item $T$, the rows of $T_0$: a list of lists of values.
\end{itemize}
A \emph{state} is a list $s$ of length $n$ with $s_i < d_i$. It satisfies the condition $(i, v)$
when $i < n$ and $s_i = v$, and it is \emph{valid} when every clause of $S$ has a condition it
satisfies. Every triple is an instance; nothing is rejected as malformed, and the semantics
settles each degenerate case. A condition on a nonexistent attribute, or with a value outside
the domain, is never satisfied; an empty clause makes every state invalid; a domain of size
zero leaves no states; a row of the wrong length or with an out-of-domain value covers no state;
and repeated rows or conditions change nothing. Then
\begin{align*}
\textit{MissingNonempty} &= \{I : \text{some valid state is not a row of } T\},\\
\textit{Complete} &= \overline{\textit{MissingNonempty}},
\end{align*}
and $\textit{missingCount}(I)$ is the number of valid states that are not rows.

\subparagraph*{A running example.}
Take two attributes with $d = [3, 2]$, the constraint $A_0 = 0 \lor A_1 = 1$, and the rows
$(0,0)$, $(0,1)$, $(1,1)$:
\[
I = \bigl([3, 2],\ [[(0,0), (1,1)]],\ [[0,0], [0,1], [1,1]]\bigr).
\]
% [rev m12]
Of the six states, four are valid: $(0,0)$, $(0,1)$, $(1,1)$ and $(2,1)$. The first three are
rows, so $M(T_0) = \{(2,1)\}$: the instance is in \textit{MissingNonempty}, it is incomplete,
and $\textit{missingCount}(I) = 1$.

% [rev m1]
\subparagraph*{Encoding and input size.}
The main encoding, \lstinline{encInstBin}, writes the triple with Complexitylib's encodings of
pairs and lists and every number in binary: the domain sizes, the attribute index and value of
each condition, and the values of each row. The number of attributes is not written separately;
it is the length of the list~$d$. The input size $|x|$ in every polynomial bound below is the
length of this encoding of the instance~$x$. With binary domain sizes,
$|\mathcal{D}| = \prod_i d_i$ can be exponential in $|x|$ even for one attribute.

% [rev m6] [rev m12]
Checking a condition $(i, v)$ against a state reads $s_i$ at a binary index~$i$. We use
Complexitylib's \lstinline{listGetD_binary}, which walks the list while reading the bits of the
index and returns a default when the list ends, so the cost is polynomial in the lengths of the
list and of the index's binary encoding; the index is never expanded to unary. The values are
shifted by one before the lookup, so that the default~$0$ matches no value: an out-of-range
attribute fails the condition, as the semantics requires. A second encoding,
\lstinline{encInst}, writes the attribute index in unary. The identity is polynomial-time from
\lstinline{encInst} to \lstinline{encInstBin} (\lstinline{encInst_to_encInstBin_polyTime}), so
hardness for the unary encoding transfers to the binary one. The NP-, coNP- and
\#P-completeness results of \Cref{sec:decision,sec:counting,sec:monotone} are proved for both
encodings, and the polynomial-time results of \Cref{sec:boundary} for \lstinline{encInstBin},
each with a \lstinline{_unary} form for \lstinline{encInst} (\Cref{app:statements}). The
hardness results on the Boolean subclass (\Cref{sec:decision}) and the explicit-degree bounds of
\Cref{sec:boundary} are stated for \lstinline{encInstBin} only.

\subparagraph*{Adequacy.}
An instance denotes a table in latc-check's model: attributes \lstinline{Fin n}, domains
\lstinline{Fin dᵢ}, the constraint predicate given by the clauses, and as $T_0$ the tuples whose
list of values is a row. The word problems are exactly latc-check's:
\begin{lstlisting}
theorem mem_MissingNonempty_iff (I : Inst) :
    I ∈ MissingNonempty ↔ (missing I.con I.table).Nonempty
theorem mem_Complete_iff (I : Inst) : I ∈ Complete ↔ missing I.con I.table = ∅
\end{lstlisting}
A third theorem, \lstinline{coe_missing_map_toList}, shows that the tuples of $M(T_0)$, read as
lists, are exactly the witnesses of \textit{MissingNonempty}, and \textit{missingCount} is
defined as the cardinality of latc-check's missing set. Every complexity result below is
therefore a statement about latc-check's \lstinline{missing} on clause-constrained tables; it
says nothing about constraint languages other than clauses.

%-----------------------------------------------------------------------------------------------
\section{The decision problems}
\label{sec:decision}

\begin{theorem}
\label{thm:decision}
Under \lstinline{encInstBin}, \textit{MissingNonempty} is NP-complete and \textit{Complete} is
coNP-complete.
\end{theorem}

\begin{proof}[Proof sketch]
\emph{Membership.} The certificate is a missing state. The verifier checks that it lies in the
domains, satisfies every clause and is not a row; it runs in polynomial time
(\lstinline{checkBin_polyTime}). A state in the domains has $s_i < d_i$ for every~$i$, so its
encoding is no longer than that of~$d$, which bounds the certificate length by~$|x|$.
\emph{Hardness.} Given a CNF formula $\varphi$, renumber its variables densely as
$0, \dots, n-1$ and build the instance with $n$ attributes of domain $\{0,1\}$, one clause per
clause of $\varphi$ (the literal $x_i$ becomes $A_i = 1$, and $\lnot x_i$ becomes $A_i = 0$), and
no rows. Its valid states are the satisfying assignments of~$\varphi$, and with no rows they are
all missing. So $\varphi$ is satisfiable iff the instance is in \textit{MissingNonempty}, and
unsatisfiable iff it is complete: a reduction from SAT and from UNSAT, polynomial-time
(\lstinline{General.fromSAT_polyTime}). Composed with Complexitylib's Cook--Levin theorems and with the
identity re-encoding, this gives the theorem (\lstinline{MissingNonempty_npComplete_bin},
\lstinline{Complete_coNPComplete_bin}). By adequacy the same holds for
$\{I : M(T_0) \neq \emptyset\}$ and $\{I : M(T_0) = \emptyset\}$ stated with latc-check's
\lstinline{missing}.
\end{proof}

% [rev D5]
\subparagraph*{Hardness on Boolean domains with empty tables.}
The reduction produces only instances whose domains are all $\{0,1\}$ and whose table is empty:
writing \textit{BoolEmpty} for the class of such instances, \lstinline{fromSAT_mem_boolEmpty}
states that every $\mathit{fromSAT}(\varphi)$ is in it. Hardness therefore holds already on
that class, and Lean states it under \lstinline{encInstBin}, with the convention of
\Cref{sec:monotone} that a restricted problem contains its restriction:
$\textit{MissingNonempty} \cap \textit{BoolEmpty}$ is NP-hard
(\lstinline{MissingNonemptyBoolEmpty_npHard_bin}) and $\textit{Complete} \cap
\textit{BoolEmpty}$ coNP-hard (\lstinline{CompleteBoolEmpty_coNPHard_bin}). The promise
version, which asks to tell the instances of \textit{BoolEmpty} with $M(T_0) \neq \emptyset$
from those with $M(T_0) = \emptyset$ and allows any answer elsewhere, is NP-hard as a promise
problem in Complexitylib's sense (\lstinline{missingNonemptyPromise_npHard_bin}). The counting
reduction of \Cref{thm:counting} is the same map, and gives the same for counting
(\lstinline{missingCountBoolEmpty_sharpPHard_bin}).

% [rev m8] [rev m12]
The empty table shows that the hardness lies in the constraints: with $T_0 = \emptyset$,
completeness is unsatisfiability of the constraints. Rows that are complete tuples only remove
single states, and \Cref{sec:boundary} shows that without constraints they cannot make the
problems hard; compact rows could (\Cref{sec:related}).

%-----------------------------------------------------------------------------------------------
\section{Counting}
\label{sec:counting}

\begin{theorem}
\label{thm:counting}
Under \lstinline{encInstBin}, \textit{missingCount} is \#P-complete under parsimonious
reductions.
\end{theorem}

\begin{proof}[Proof sketch]
\emph{Membership.} Membership needs every missing state to have exactly one accepted
certificate. Let $L$ be the length of the encoding of the domain sizes~$d$. A certificate is a
bit string of length~$L$. The verifier decodes the first $n$ numbers from it into a state~$s$
and accepts when $s$ is missing \emph{and} the certificate equals the canonical certificate of
$s$: the encoding of $s$ followed by zeros up to length~$L$. The encoding of a state in the
domains is no longer than that of $d$, so every missing state has a canonical certificate of
length~$L$; the canonical certificate determines the state; and any other string decoding to $s$
is rejected. So the accepted certificates of length $L$ are in bijection with $M(T_0)$
(\lstinline{missingCount_eq_acceptCount}), and both the verifier and $L$ are polynomial-time.
Without the canonicity test, a state could have several accepted certificates (different
padding) and the count would be wrong.

\emph{Hardness.} The reduction of \Cref{thm:decision} is parsimonious: the satisfying
assignments of $\varphi$ and the missing states of its instance are in bijection, so
$\#\mathrm{SAT}(\varphi) = \textit{missingCount}(\mathit{fromSAT}(\varphi))$
(\lstinline{sharpSAT_eq_missingCount}). Composition with Complexitylib's parsimonious
Cook--Levin theorem gives \lstinline{missingCount_sharpPComplete_bin}.
\end{proof}

%-----------------------------------------------------------------------------------------------
\section{Monotone two-literal constraints and Valiant's theorem}
\label{sec:monotone}

% [rev D5]
Russell~\cite{PaperA} notes that counting the valid states stays \#P-complete when every
clause is $A_i = 1 \lor A_j = 1$ over domains $\{0,1\}$, citing
Valiant~\cite{Valiant1979enumeration}. Deciding such instances is easy: with domains
$\{0,1\}$ and only positive two-literal clauses on existing attributes, the all-ones state is
valid; the shape tests also require every condition to name an existing attribute, so every
such instance with no rows has a missing state (\lstinline{mem_missingNonempty_of_pos2B}). With
rows, deciding is still polynomial\informal: a partial assignment to a positive two-literal CNF
extends to a model if and only if no clause has both literals set to~$0$ (set the rest to~$1$),
so backtracking lists the valid states with polynomial delay, and listing $|T|+1$ of them
decides whether one is missing. So neither count can be \#P-hard under parsimonious reductions
unless P~=~NP\informal, and the right notion is the one-call reduction of \Cref{def:onecall}.

% [rev D3]
\subparagraph*{Restricted problems.}
A problem restricted to a class $K$ of instances is stated in this development as a problem on
all instances that contains its restriction. A restricted decision problem is the set
$\{I : I \in K \land I \in A\}$, so that its decider also checks $I \in K$, and a restricted
count returns $0$ outside~$K$, as Complexitylib's \#3SAT does. The counts of this section follow
the convention, and so do the restricted sets of \Cref{sec:decision,sec:boundary}. It is
harmless because every class used is decidable in polynomial time: the shape tests below
(\lstinline{pos2B_polyTime_bin}, \lstinline{pos2AnyB_polyTime_bin}), \textit{BoolEmpty}, ``no
constraints'' and ``at most $k$ attributes'', the last three plain comparisons on the input. The
\emph{promise version} of a restricted problem asks for the right answer only on the instances
of~$K$. A polynomial-time decider for the restricted set solves the promise version, so the
upper bounds of \Cref{sec:boundary} hold for both forms; conversely, a solver for the promise
version, run after the test for~$K$, decides the restricted set\informal. The hardness proofs
map their inputs into~$K$, so they show the promise versions hard as well. Lean states promise
versions in two places: \lstinline{missingNonemptyPromise_npHard_bin} (\Cref{sec:decision}),
and Statement~2 of the explicit degree (\Cref{sec:boundary}), whose input type is the subtype of
instances with at most $k$ attributes. For the counts below, the promise version is the count
on the instances that pass the shape test; the reduction leaves that class only on formulas
where both sides are~$0$ by the convention, and the promise versions of the counts are not
stated separately in Lean\informal.

% [rev D3]
Let \textit{SharpMonotone2SAT} be the number of models of a CNF formula whose clauses are all
$x_i \lor x_j$, and \textit{countValidMonotone2} the number of valid states of an instance with
domains $\{0,1\}$, clauses of two conditions with value~1 on existing attributes, and no rows. Let
\textit{countMissingMonotone2} be the same without the condition on rows: the number of missing
combinations $|M(T_0)|$ of an instance with those domains and constraints and any rows
(\lstinline{countMissingMonotone2_eq_card_missing}). As with Complexitylib's \#3SAT, the three
functions are total and return~$0$ on inputs of any other shape. Both counts are in \#P, and
\#Monotone-2SAT reduces to each parsimoniously (\lstinline{sharpMonotone2SAT_reduces}, and
\lstinline{sharpMonotone2SAT_reduces_countMissing} in \lstinline{Latc.PaperD.Valiant}): the
reduction produces either a monotone instance with no rows, where the two counts agree, or a
fixed instance of the wrong shape, where both are~$0$. What remains is the hardness of
\#Monotone-2SAT itself.

\begin{theorem}[one-call Valiant]
\label{thm:valiant}
\#Monotone-2SAT is one-call \#P-hard. Hence \textit{countValidMonotone2} and
\textit{countMissingMonotone2} are one-call \#P-complete under \lstinline{encInstBin} and
\lstinline{encInst}.
\end{theorem}

% [rev D7] [rev m2]
Valiant proved \#Monotone-2SAT \#P-complete under polynomial-time Turing
reductions~\cite{Valiant1979enumeration}. Its relation $x \lor y$ is not affine, so it also lies
on the hard side of Creignou and Hermann's dichotomy for counting generalised
satisfiability~\cite{CreignouHermann1996}. Senellart and
Gnatenko's Lean library~\cite{DescriptiveComplexityLib} machine-checks a one-call form in its
descriptive-complexity setting, by a chain of four reductions: from \#1-in-SAT to counting
cliques, to counting independent sets of a given size, to counting all independent sets with one
call, by attaching pendant leaves to every vertex and reading one digit of the answer, and to
\#Monotone-2SAT. Our proof uses the same two devices, a parsimonious encoding of models as
independent sets of a fixed size, and pendant-leaf amplification followed by digit extraction.
It differs from theirs in fusing the chain into one reduction from
\#SAT, through the ``first true literal'' graph below, which handles any CNF and needs no
counting problem on graphs. What it adds is a new formal proof in Complexitylib's model, where
every step is a polynomial-time Turing-machine computation; it claims no new mathematics. Since
\#SAT is \#P-hard under parsimonious reductions, by \Cref{def:onecall} it suffices to reduce
\#SAT to \#Monotone-2SAT with one call.

\subparagraph*{Step 1: models as independent sets.}
% [rev m1] (independent sets are J)
Let $\varphi$ have variables $0, \dots, n-1$ (after dense renumbering) and clauses
$c_1, \dots, c_m$. Build a graph $G$ with a vertex $(i, b)$ for each variable $i$ and value $b$,
and a vertex $o$ for each literal occurrence, read ``$o$ is the first true literal of its
clause''. Edges join $(i, 0)$ and $(i, 1)$; any two occurrences of the same clause; an
occurrence and the value of its variable that falsifies it; and an occurrence and each value
that makes an \emph{earlier} literal of its clause true. The vertices fall into $K = n + m$
cliques, $\{(i,0),(i,1)\}$ for each variable and the occurrences of each clause, so every
independent set has at most $K$ vertices. An independent set of exactly $K$ vertices picks a
value for each variable, which is an assignment~$\sigma$, and an occurrence in each clause that
$\sigma$ makes true, with every earlier literal false. That occurrence is unique given~$\sigma$,
so these independent sets are in bijection with the models of~$\varphi$
(\lstinline{modelCount_eq_card}). The ``first true literal'' choice is what makes the
correspondence exact, for any CNF with densely numbered variables (the hypotheses of
\lstinline{modelCount_eq_card}): repeated and complementary literals need no special case. An
empty clause has no vertex, so no independent set reaches $K$, and the count is $0$, as it should
be.

For example, $\varphi = x_0 \lor x_1$ has $n = 2$, $m = 1$, $K = 3$, vertices $(0,0)$,
$(0,1)$, $(1,0)$, $(1,1)$, $o_1$ for $x_0$, $o_2$ for $x_1$. The edges at $o_1$ go to $o_2$ and
$(0,0)$; those at $o_2$ go to $o_1$, $(1,0)$ and $(0,1)$. The independent sets of size~3 are
$\{(0,1), (1,0), o_1\}$, $\{(0,1), (1,1), o_1\}$ and $\{(0,0), (1,1), o_2\}$: three, one per model.

\subparagraph*{Step 2: a monotone formula.}
% [rev m1] (the number of vertices is \nu, the oracle answer y)
Let $\nu$ be the number of vertices of $G$, and give every vertex $t = n + 1$ fresh pendant
leaves. The monotone formula $\psi$ has a variable per vertex and leaf, and a clause $x_u \lor
x_w$ per edge $\{u, w\}$ of the graph with leaves. Its models are the vertex covers, the
complements of the independent sets. An independent set of the graph with leaves consists of an
independent set $J$ of $G$ together with any subset of the leaves of the $\nu - |J|$ vertices
outside~$J$, so
\[
y \;=\; \#\mathrm{SAT}(\psi) \;=\; \sum_{J \text{ independent in } G} 2^{\,t\,(\nu - |J|)}
\qquad (\text{\lstinline{modelCount_psi}}).
\]
Every vertex has leaves, so every variable of $\psi$ occurs in some clause, and $\psi$ has exactly two
positive literals per clause.

% [rev m1]
\subparagraph*{Step 3: one digit.}
Write $y$ in base $2^t$. Multiply by $2^{tK}$ and divide by $2^{t\nu}$:
\[
\Bigl\lfloor \frac{y \cdot 2^{tK}}{2^{t\nu}} \Bigr\rfloor
= \Bigl\lfloor \sum_{J} 2^{\,t\,(K - |J|)} \Bigr\rfloor .
\]
Since $|J| \le K$, every exponent is a nonnegative multiple of $t$, so the sum is an integer and
the floor is exact. Sets with $|J| < K$ contribute multiples of $2^t$, and sets with $|J| = K$
contribute $1$ each. Their number is $\#\mathrm{SAT}(\varphi) \le 2^n < 2^t$. So
\[
\#\mathrm{SAT}(\varphi) \;=\; \Bigl\lfloor \frac{y \cdot 2^{tK}}{2^{t\nu}} \Bigr\rfloor \bmod 2^t
\qquad (\text{\lstinline{sharpSAT_eq_post}}).
\]
Multiplying before dividing avoids computing $\nu - K$, which can be negative when $\varphi$ has
empty clauses ($K$ counts them, $\nu$ does not). Then every term $2^{t(K-|J|)}$ has $|J| < K$, so the
quotient is a multiple of $2^t$ and the digit is $0$, which is correct.
The post-processing needs only binary multiplication, division, remainder and powers of two on
numbers of polynomially many bits, all polynomial-time in Complexitylib; the construction of
$\psi$ is polynomial-time; and the oracle is called once, on~$\psi$
(\lstinline{toMonotone_polyTime}, \lstinline{post_polyTime}, \lstinline{sharpSAT_oneCallReducible}).
Composing with the parsimonious \#P-hardness of \#SAT gives
\lstinline{sharpMonotone2SAT_oneCallSharpPHard}, and with the reductions to
\textit{countValidMonotone2} and \textit{countMissingMonotone2}, the second half of
\Cref{thm:valiant} (\Cref{app:statements}). The development also keeps
the conditional form, \lstinline{countValidMonotone2_oneCallSharpPComplete_bin}, which takes
\#Monotone-2SAT's hardness as a hypothesis; \Cref{thm:valiant} discharges it.

%-----------------------------------------------------------------------------------------------
\section{Where the hardness comes from}
\label{sec:boundary}

% [rev D3]
The reductions of \Cref{sec:decision,sec:counting} use clause constraints and as many
attributes as the formula has variables. Removing either ingredient makes all three problems
polynomial-time (\Cref{tab:boundary}). As in \Cref{sec:monotone}, each restricted problem
below contains its restriction.

\subparagraph*{No constraints.}
Then every state is valid and $|M(T_0)| = \prod_i d_i - r$, where $r$ is the number of
\emph{distinct} rows that are states (\lstinline{missingCount_eq_of_noConstraints}). The product
can be exponential in $|x|$, but its binary encoding is no longer than those of the factors
together, so a fold that multiplies in binary stays polynomial: each intermediate product has
at most $|x|$ bits. The count is therefore polynomial-time, and so are completeness and
non-emptiness, which compare $\prod_i d_i$ with~$r$ (\lstinline{missingCountNoConstraints_polyTime},
\lstinline{completeNoConstraints_mem_P}, \lstinline{missingNonemptyNoConstraints_mem_P}).

\subparagraph*{At most $k$ attributes: deciding.}
With binary domain sizes, $\mathcal{D}$ cannot be enumerated even for $k = 1$. Let $N$ be the
number of value entries written in the instance, in conditions and in rows, and call a number
\emph{unwritten} if it occurs in none of them. An unwritten value satisfies no condition, and a
state containing one is not a row; so in a missing state, an unwritten value can be replaced by
any other unwritten value of the same domain, and the state stays missing. At most $N$ of the
numbers $0, \dots, N$ are written, so one of them, $u$, is unwritten. If a missing state uses an
unwritten value $v > N$ in attribute~$i$, then $d_i > v > u$, so $u$ is in the domain and can
replace~$v$. Hence if a missing state exists, one exists whose values lie among $0, \dots, N$
and the written values: at most $2N+1$ candidates per attribute and $(2N+1)^k$ candidate
states (\lstinline{length_candidateStates_le},
\lstinline{exists_missing_candidate_iff}). Checking each candidate is polynomial, so both decision
problems are in P for each~$k$ (\lstinline{missingNonemptyFixedAttrs_mem_P},
\lstinline{completeFixedAttrs_mem_P}).

\subparagraph*{At most $k$ attributes: counting.}
% [rev m1] (class vectors are \gamma)
A candidate state stands for all the states obtained by replacing its unwritten values, so
counting must weigh it. For each attribute, split its domain into classes: each written value
inside the domain alone, and one class holding all the unwritten values of the domain (possibly
empty), represented by~$u$. Which
conditions a state satisfies, and whether it is a row, depend only on its vector of classes, so
\[
|M(T_0)| = \sum_{\substack{\gamma \text{ a class vector,}\\ \text{valid and not a row}}}\
\prod_{i} |\gamma_i|
\qquad (\text{\lstinline{missingCount_eq_weighted_of_le}}),
\]
over at most $(N+1)^k$ class vectors (\lstinline{length_classStates_le}). Each product is at most
$\prod_i d_i$, so the sum is computed in binary by the bounded-growth fold of the
no-constraint case, and the count is polynomial-time for each~$k$
(\lstinline{missingCountFixedAttrs_polyTime}).

\subparagraph*{An explicit degree.}
% [rev m1] [rev m4] [rev m12]
``Polynomial time for each $k$'' allows the polynomial to depend on $k$ arbitrarily. We bound
its form: the exponent is linear in the number of attributes, with constants that are proved to
exist and are not computed. The constants are written $c$, $e$, $e_0$ and~$e_1$ below; the Lean
statements call the exponent constants \lstinline{d}, \lstinline{d₀} and \lstinline{d₁}. Two
names of bound frameworks occur. The statements use \lstinline{ComputableTidilyInTimeAndSpace}
\lstinline{f encIn encOut t s} (from CSLib), which says that one machine computes $f$ on every input $a$
within exactly the time $t(a)$ and the space $s(a)$. The proofs use Complexitylib's
\lstinline{ComputableTidilyInO} \lstinline{f encIn encOut T S}, which is defined from it: for some
constant~$c$, the bounds $c\,(T(a)+1)$ and $c\,(S(a)+1)$ hold and the output has length at most
$c\,(S(a)+1)$. The second form keeps explicit bound functions and a single constant through
composition, case splits and folds; the lemma \lstinline{timeAndSpace_of_inO_xp} converts it to
the first at the end (\Cref{sec:engineering}). The following are proved, each for
non-emptiness, completeness and counting, each for time and space, under
\lstinline{encInstBin}:
\begin{enumerate}
\item \emph{One machine for all instances.} There are $c, e$ such that one machine decides
  non-emptiness on every instance within $c\,(|x|+1)^{en+e}$, where $n$ is the instance's number
  of attributes (\lstinline{missingNonempty_explicitDegree}; \Cref{app:statements}).
% [rev m5]
\item \emph{Uniform constants in $k$} (\lstinline{missingNonemptyFixedAttrs_uniform}). There are
  $c, e_0, e_1$ such that for every $k$ some machine decides the problem within
  $c\,(|x|+1)^{e_0 k + e_1}$ on every instance with at most $k$ attributes. Its inputs are the
  instances with at most $k$ attributes, encoded by \lstinline{encInstBin}: the Lean input type
  is the subtype \lstinline|{I : Inst // I.n ≤ k}|, so nothing is claimed about other inputs,
  and the statement is a promise version in the sense of \Cref{sec:monotone}. The proof takes the
  machine of~1 for every~$k$; the statement asserts only that a machine exists for each~$k$.
\item \emph{$k$ in the input.} There are $c, e$ such that one machine, given $k$ in unary and an
  instance, decides ``at most $k$ attributes and $M(T_0) \neq \emptyset$'' within
  $c\,(|(k,x)|+1)^{ek+e}$ (\lstinline{missingNonemptyWithBound_explicitDegree}).
\end{enumerate}
% [rev m7] [rev D5] [rev m12]
The degree grows linearly with the number of attributes. This is an XP bound, $|x|^{O(k)}$;
fixed-parameter tractability would need $f(k) \cdot |x|^{O(1)}$, and it is unlikely\informal.
With $k$ attributes, non-emptiness is a constraint satisfaction problem with $k$ variables,
and such problems are W[1]-hard with the number of variables as parameter: Papadimitriou and
Yannakakis prove it for conjunctive queries, by a reduction from Clique~\cite{PapadimitriouYannakakis1999}, a result already stated as Theorem~1 of its
conference version (PODS 1997). An argument we have not formalised reduces Multicolored Clique, which is
W[1]-complete~\cite[Theorem~13.25]{CyganEtAl2015}, to non-emptiness with $k$ attributes
directly\informal. Give attribute $i$ the vertices of colour class $i$ as values, no rows, and
for each non-adjacent pair $u, w$ in classes $i \ne j$ the clause
$\bigvee_{v \ne u} A_i = v \lor \bigvee_{v \ne w} A_j = v$ (``not both $A_i = u$ and
$A_j = w$''); the valid states are exactly the multicoloured cliques.

% [rev m12]
\emph{What is not proved.} For each fixed $k$, a separate machine with $k$ built in that rejects
instances with more than $k$ attributes, with constants independent of~$k$, is not proved.
Testing $n \le k$ inside the machine goes through Complexitylib's case-split and
constant-output lemmas, whose constants depend on the machine, and making them uniform would
need changes to the library. Statements~2 and~3 cover the same ground: the first on the
instances where the answer can be positive, the second on all instances with $k$ given as
input.

% [rev D5]
\begin{table}[t]
\caption{Where the hardness comes from. A dagger marks a claim argued on paper only; every
other entry is a Lean theorem. The hardness of the first row holds already on Boolean domains
with $T_0 = \emptyset$.}
\label{tab:boundary}
\centering\small
\setlength{\tabcolsep}{5pt}% layout only
\begin{tabular}{L{0.11\textwidth}L{0.13\textwidth}L{0.12\textwidth}L{0.15\textwidth}L{0.16\textwidth}L{0.14\textwidth}}
\hline
rows & constraints & attributes & deciding & counting & hardness via \\
\hline
tuples & clauses & unbounded & NP-, coNP-complete & \#P-complete & many-one, parsimonious \\
tuples & none & unbounded & P & polynomial-time & -- \\
tuples & clauses & at most $k$ & P, degree $O(k)$ & polynomial-time, degree $O(k)$ & -- \\
tuples & monotone two-literal & unbounded & P\informal{} (trivial if $T_0 = \emptyset$) & one-call \#P-complete & one-call \\
compact (wildcards) & none & unbounded & coNP-complete\informal & \#P-complete\informal & DNF tautology, \#SAT \\
\hline
\end{tabular}
\end{table}

%-----------------------------------------------------------------------------------------------
\section{Related work}
\label{sec:related}

\subparagraph*{Compact rows.}
% [rev D5] [rev m11] [rev m12] [rev m2]
Our rows are complete tuples, as in latc-check. In DMN and in many rule bases a row is compact: a
conjunction of tests over the declared finite domains, with wildcards or ranges, standing for
many tuples. Then the rows form a DNF formula and the uncovered inputs are the models of its
negation, a CNF. Both directions hold\informal: any CNF $\varphi$ arises this way, from the
wildcard rows given by the terms of $\neg\varphi$, with the uncovered tuples exactly the models
of~$\varphi$. So, even without constraints, completeness is the DNF tautology problem,
coNP-complete, and counting the uncovered inputs is parsimoniously equivalent to \#SAT, hence
\#P-complete by the parsimonious Cook--Levin theorem; tautology of DNF is coNP-complete by Cook's
theorem~\cite{Cook1971}; membership holds because an uncovered tuple is a certificate of
polynomial size, even with domain bounds written in binary, checked against each row in
polynomial time. Russell~\cite{PaperA} makes the same observation. Elmallah and Gouda's firewall rules are compact rows of this kind: each rule tests every field
against an interval of integers~\cite{ElmallahGouda2017}.
Thiemann and Yamada's coNP-hardness reduction turns each clause into a row of true, false and
wildcard patterns, so it already shows coNP-hardness for decision tables with wildcard
rows~\cite{ThiemannYamada2026}; what this paper adds is the case of rows that are complete
tuples, constraints, counting and the mechanisation. With complete tuples as rows, as here, the
hardness comes from the constraints (\Cref{sec:boundary}).

\subparagraph*{Hardness of coverage.}
% [rev m2] [rev m11]
Hardness of rule coverage is not new in itself. Elmallah and Gouda prove firewall completeness NP-hard, by a reduction from a
slice-probing problem that they reduce from 3-SAT~\cite{ElmallahGouda2017}. The reduction maps
to the complement of completeness, so by its form it shows coNP-hardness; for membership they
only remark that some of their problems can be shown to be in NP and the others in coNP. Our account follows their 2013 technical report of the same title.
Thiemann
and Yamada prove pattern completeness coNP-hard and verify their decision procedure in
Isabelle, the hardness proof being on paper~\cite{ThiemannYamada2026}. According to its abstract, the paper of Kapur et al.\ shows sufficient completeness coNP-complete for term rewriting systems with
free constructors, and state that their complexity results also apply to ground
reducibility~\cite{KapurNRZ1991}; Thiemann and Yamada cite the paper for the coNP-completeness of
quasi-reducibility.
Deciding whether a point lies outside a union of boxes is NP-hard~\cite{Chan2010}, and computing
the volume of the union is \#P-hard~\cite{BringmannFriedrich2010}. For firewalls with a fixed
number of header fields, Hamilton, Roughan and Nguyen give polynomial-time
algorithms~\cite{HamiltonRoughanNguyen2022}, an analogue of our fixed-$k$ results.

\subparagraph*{Decision tables and tabular specifications.}
% [rev m12]
Checking tables for completeness and consistency is long established: in expert-system
checkers~\cite{SuwaScottShortliffe1982}, for decision tables~\cite{VanthienenMA1998}, for SCR and
RSML requirements~\cite{HeitmeyerJL1996,HeimdahlLeveson1996}, and for DMN, where Calvanese et al.\
give analysis algorithms~\cite{CalvaneseDLMMT2016,CalvaneseDLMMT2018} and, with background
knowledge, an ExpTime upper bound~\cite{CalvaneseMDM2019}. PVS generates coverage and
disjointness obligations for tabular specifications~\cite{OwreRushbyShankar1997}; these are
proof obligations about given tables, and come with no complexity results. Menghi et al.\ prove
completeness of Simulink Requirements Tables undecidable, by reduction from Hilbert's tenth
problem, because their conditions are arithmetic over infinite domains, and check it with
SMT~\cite{MenghiEtAl2025,ZanengaEtAl2025}. Vu and Demetrovics study decision tables whose entries
may be unknown: for one given input, deciding its possible and certain decisions is NP- and
coNP-complete, and counting the valuations yielding a decision \#P-hard, all by encoding a formula
in the constraints on the unknowns; a fixed number of unknowns is polynomial, another analogue of
our fixed-$k$ case~\cite{VuDemetrovics2026}. They define completeness for a single input but state
no complexity for deciding it. Their model asks what a partly unknown table says about one input;
ours asks whether a fully specified table covers every input its declared domains and constraints
allow.

\subparagraph*{Rule-base verification.}
Levy and Rousset gave the first complexity results for verifying rule
bases~\cite{LevyRousset1996,LevyRousset1998}: whether no legal input derives a bad output, for
relational inputs and Horn rules, NP-complete to undecidable depending on recursion, and
$\Pi_2^p$-complete with order predicates and negation. Their ``input completeness'' constraints
restrict which inputs are legal, and they leave output completeness open. Coverage of a
propositional table could be phrased in their framework, but it lies in coNP~$\subseteq \Pi_2^p$:
coNP-complete for completeness, as shown here, with NP for non-emptiness and \#P for counting.

Apart from the argument on paper in~\cite{PaperA}, and to our knowledge, none of these works
classifies the completeness of decision tables with complete
tuples as rows over declared finite domains, or counts their missing combinations, and none of the
hardness proofs is machine-checked.

\subparagraph*{Mechanised computability and complexity.}
% [rev m13] [rev m12]
Cook--Levin has been mechanised in ACL2~\cite{GamboaCowles2004}, Coq~\cite{GaeherKunze2021} and
Isabelle~\cite{Cook_Levin-AFP}; reductions between Karp's problems are verified in
Isabelle~\cite{PolyReductionsIsabelle}; and Senellart and Gnatenko prove all 21 of Karp's
problems NP-complete in Lean~\cite{SenellartGnatenko2026}, with NP defined in
descriptive-complexity terms and first-order reductions; their library also proves
\#Monotone-2SAT one-call \#P-complete in that setting~\cite{DescriptiveComplexityLib}. Mathlib~\cite{mathlib2020} itself defines polynomial-time
computability for its stack-based Turing machines: \lstinline{TM2ComputableInPolyTime}, in the
file \texttt{Mathlib/Computability/TuringMachine/Computable.lean} by Spelier and van
Gent~\cite{MathlibTMComputable}. Turing machines
and computability have been mechanised in Isabelle/HOL~\cite{XuZhangUrban2013} and Coq~\cite{ForsterKunzeWuttke2020}, computability through
partial recursive functions in Lean~\cite{Carneiro2019}, and time complexity through the weak
call-by-value $\lambda$-calculus in Coq~\cite{ForsterKSW2021}. Our setting needs explicit
encodings and running times; in exchange its classes are the textbook ones on Turing machines,
so polynomial time, binary encodings of numbers and the size of the input, the quantities the
tractability results of \Cref{sec:boundary} are about, are first-class.

%-----------------------------------------------------------------------------------------------
% [rev D1]
\section{Trusted base and audit}
\label{sec:formal}

\subparagraph*{Trust.}
% [rev D2] [rev D5]
The axiom audit prints the axioms of 328 of the development's 345 theorem and lemma
declarations, among them every theorem this paper cites (the other 17 are helper lemmas used by
audited theorems); each uses only \lstinline{propext}, \lstinline{Classical.choice} and
\lstinline{Quot.sound}, or fewer, and none uses \lstinline{sorryAx}. This shows that no proof is
left unfinished and no axiom is added. It does not show that the definitions are the intended
ones. The trusted base is Lean's kernel (we did not replay the proofs in an external checker);
Complexitylib's definitions of machines, running time and space
(\lstinline{ComputableTidilyInTimeAndSpace}, which carries the explicit-degree claims),
encodings, the classes P, NP, coNP and \#P, the reductions, promise problems and the notions of
completeness built on them, and its CNF semantics and \#SAT; our encodings of instances;
\Cref{def:onecall}; and the problem definitions themselves: \textit{MissingNonempty} with
\lstinline{InDomains}, \lstinline{CondHolds} and \lstinline{Valid}, \textit{missingCount}, the
monotone shape tests, \textit{SharpMonotone2SAT} and \textit{BoolEmpty}. \Cref{app:definitions}
copies every one of our trusted definitions verbatim from the source. The adequacy theorems
shrink the list: the general problems are proved equal to latc-check's, so for them a reader need
trust only latc-check's definition of the missing set and our map from an instance to the table
it denotes (\lstinline{Inst.con}, \lstinline{Inst.table}), each a few lines
(\Cref{sec:prelim,sec:encoding}). For the monotone counts, adequacy holds only on instances that
pass the shape test, so the shape test is trusted; an earlier version of it accepted conditions on
nonexistent attributes, which the present one rejects.

\subparagraph*{Worst case and practice.}
% [rev m12]
The classification rules out a polynomial algorithm for all tables with clause constraints,
unless P~=~NP; it does not say that tables met in practice are slow to check. For tables whose
rows are complete tuples, few attributes or the absence of constraints give the polynomial
regimes of \Cref{sec:boundary}, and the difficulty comes from constraints across many
attributes; neither the size of the domains nor that of the table makes the problems hard. For
compact rows, as in DMN, the rows themselves can carry the hardness (\Cref{sec:related}), and
our tractability results do not apply.

%-----------------------------------------------------------------------------------------------
% [rev D1] New section, replacing the paragraph "Engineering" of the former Section 9.
\section{Engineering the development}
\label{sec:engineering}

This section reports what the formalisation took and what it taught. Its figures come from the
repository at the commit described in \Cref{app:artifact} and from the commit history, which is
summarised in the repository's README. No working hours were recorded, so it reports lines,
commits and dates only.

\subsection{Size and effort per part}
\label{sec:effort}

% [rev D1]
\Cref{tab:modules} lists the modules. The development has 4\,520 lines of Lean in 15 files,
besides the two axiom-audit files (383 lines), with 345 theorem and lemma declarations; it sits
on latc-check's model, copied unchanged into \texttt{LatcFormal/}. By lines, the effort fell
unevenly on four parts:
\begin{itemize}
\item \emph{Encodings and the decision problems}, 411 lines (\texttt{General.lean} with the unary
  index, \texttt{Binary.lean} with the binary one). Most of it is the certificate check and its
  polynomial time; the binary index cost 206 lines, mainly the re-encoding and the shifted lookup
  of \Cref{sec:encoding}.
\item \emph{Adequacy}, 152 lines (\texttt{Adequacy.lean}): the map from a state to a tuple of
  \lstinline{Fin} values, its injectivity and the two adequacy theorems.
\item \emph{The explicit degree}, 692 lines (\texttt{ExplicitDegree.lean}), for a statement that
  is a corollary on paper: almost as much as the 825 lines of encodings, adequacy and counting
  together (\texttt{General}, \texttt{Binary}, \texttt{Adequacy}, \texttt{CountingGeneral}).
  The cost is bookkeeping: every composition step must carry an explicit bound
  (\Cref{sec:instance-bounds}).
\item \emph{Valiant's theorem}, 1\,140 lines in four files, the largest part; 518 of them
  (\texttt{Valiant/Counting.lean}) prove the bijection between models and independent sets of
  size~$K$.
\end{itemize}
The tractable side (\texttt{Tractable.lean}, 1\,099 lines) is the largest single file: the
enumeration of candidate states needs an induction on $k$ for its polynomial time, and the
counting with classes adds the partition of each domain. The repository records commits, not
hours. Its first Lean commit (28 October 2026) holds the Boolean model and the general NP- and
coNP-completeness; adequacy, the two counting files and the conditional monotone case followed
the same day, the binary index and the tractable side on 29 October, and the three remaining
results on 31 October, after a written plan (\texttt{docs/gaps-plan.md}): counting with at most
$k$ attributes, Valiant's theorem and the explicit degree, committed 14, 31 and 35 minutes after
the plan. Commit times bound neither when the work on a part began nor the author's time
reviewing it. The three parts were developed in parallel sessions.

% [rev D1]
\begin{table}[t]
\caption{The modules of \texttt{LatcComplexity/}: lines, theorem and lemma declarations, and
definitions (\texttt{def}, \texttt{abbrev}, \texttt{instance}). The two audit files
are not counted in the total.}
\label{tab:modules}
\centering\small
\begin{tabular}{L{0.27\textwidth}R{0.07\textwidth}R{0.08\textwidth}R{0.06\textwidth}L{0.36\textwidth}}
\hline
file & lines & thms & defs & contents \\
\hline
\texttt{Bridge.lean} & 160 & 16 & 11 & Boolean model; agreement with latc-check's clauses \\
\texttt{General.lean} & 205 & 15 & 13 & instances, \texttt{encInst}, NP- and coNP-completeness \\
\texttt{Adequacy.lean} & 152 & 13 & 6 & the table an instance denotes; adequacy \\
\texttt{Counting.lean} & 140 & 10 & 2 & Boolean model: counting in \#P \\
\texttt{CountingGeneral.lean} & 262 & 20 & 7 & \textit{missingCount}: \#P-completeness \\
\texttt{Monotone.lean} & 238 & 18 & 10 & one-call reductions; monotone shape, conditional case \\
\texttt{Binary.lean} & 206 & 22 & 3 & \texttt{encInstBin}; every result above for it \\
\texttt{Tractable.lean} & 1\,099 & 81 & 26 & no constraints; at most $k$ attributes \\
\texttt{ExplicitDegree.lean} & 692 & 46 & 7 & explicit degree, uniform in $k$ \\
\texttt{Restricted.lean} & 102 & 10 & 6 & hardness on Boolean domains, $T_0 = \emptyset$ \\
\texttt{Valiant/Gadget.lean} & 227 & 19 & 14 & the graph, $\psi$, \textit{post}; polynomial time \\
\texttt{Valiant/Counting.lean} & 518 & 38 & 7 & models and independent sets of size $K$ \\
\texttt{Valiant/Pendant.lean} & 195 & 9 & 0 & the model count of $\psi$ \\
\texttt{Valiant/OneCall.lean} & 200 & 14 & 0 & the digit; one-call Valiant \\
\texttt{Valiant/AnyTable.lean} & 124 & 14 & 2 & \textit{countMissingMonotone2} \\
\hline
total & 4\,520 & 345 & 114 & \\
\hline
\end{tabular}
\end{table}

\subsection{What Complexitylib provided and what it lacked}
\label{sec:gaps}

% [rev D1]
Complexitylib provided the machine model, the encodings, the classes, the Cook--Levin theorem
with its parsimonious form, and a large stock of polynomial-time lemmas, among them binary
arithmetic (multiplication, division, remainder, powers of two), list folds, the binary list
lookup \lstinline{listGetD_binary}, a fold whose intermediate results grow boundedly
(\lstinline{foldlWithInput_of_bounded_growth}, used for the binary product of the domain sizes),
the padding lemma \lstinline{mem_SharpP_of_polyTime_length}, the bound framework
\lstinline{ComputableTidilyInO} and promise problems. The development needed five things the
library did not have, and worked around each:
\begin{enumerate}
\item \emph{Reductions with one oracle call.} The library has many-one, parsimonious and
  Turing reductions. \Cref{def:onecall} is defined in \texttt{Monotone.lean}, with its
  composition with parsimonious reductions on either side
  (\lstinline{ParsimoniousReducible.trans_oneCall},
  \lstinline{OneCallSharpPHard.trans_parsimonious}).
\item \emph{\#P-hardness beyond the SAT family.} The library has \#SAT and \#3SAT, a decision
  problem for independent sets, and no counting problem on graphs. The plan was to port
  Senellart and Gnatenko's chain of three reductions through counting problems on graphs; the
  proof instead builds the graph of \Cref{sec:monotone} directly as a CNF over numbered vertices
  and needs no graph library.
\item \emph{Explicit bounds.} The composition lemmas for polynomial time are existential: a
  composed bound exists, but its degree is not visible. The explicit degree goes through
  \lstinline{ComputableTidilyInO} instead (\Cref{sec:instance-bounds}).
\item \emph{Uniform constants.} The constants of the case-split and constant-output lemmas depend
  on the machines involved, which is why the per-$k$ machine with constants independent of~$k$
  is not proved (\Cref{sec:boundary}).
\item \emph{Reading a list at a binary index.} The first encoding wrote the attribute index in
  unary, because the list lookup used for the certificate check
  (\lstinline{ComputableTidilyInPolyTime.getElem?}) takes a unary position. The binary encoding
  then used the library's \lstinline{listGetD_binary}, with the values shifted by one so that
  the default of an out-of-range lookup matches no value.
\end{enumerate}
The library changes daily and sits on a release-candidate toolchain, so the development pins
one commit and one toolchain (\Cref{app:artifact}).

\subsection{Polynomial time by automation}
\label{sec:funprop}

% [rev D1]
Complexitylib registers its polynomial-time lemmas with Mathlib's \lstinline{fun_prop} tactic,
which proves a polynomial-time claim for a term built from registered functions by composing
their lemmas. The library's documentation asks that local definitions be unfolded so that
\lstinline{fun_prop} can see their bodies. Most polynomial-time proofs of the development are
then three lines, as for the reduction from SAT (\Cref{lst:fromsat}). The \lstinline{tidy}
forms are machines that restore their tapes, which composition needs; \lstinline{.of_tidy}
returns to the plain notion.

\begin{lstlisting}[float=t,caption={The polynomial time of the SAT reduction (\texttt{General.lean}). The two \texttt{have} lines bring library lemmas for the dense relabelling and the number of variables into the context; \texttt{fun\_prop} composes them with the lemmas for pairs, lists, \texttt{replicate}, \texttt{map} and \texttt{if}.},label={lst:fromsat}]
theorem fromSAT_polyTime : ComputableInPolyTime fromSAT encCNF encInst := by
  have h1 := (CNF.computableInPolyTime_denseRelabel_unary encNat).tidy
  have h2 := SAT.computableInPolyTime_numVars.tidy
  refine .of_tidy ?_
  unfold fromSAT toCond
  fun_prop
\end{lstlisting}

% [rev D1]
Unfolding everything does not scale. For the graph of \Cref{sec:monotone}, whose vertices are
numbered by Cantor pairing ($3\cdot\mathrm{pair}(i,b)$ for the value vertices, and similar codes
for occurrences and leaves), the remedy was to give each helper, the three vertex codes and the edge, a lemma
in two forms: one for the helper itself, proved by unfolding, and a composable one, registered
with \lstinline{fun_prop}, for the helper applied to any polynomial-time arguments
(\Cref{lst:vv}). \lstinline{fun_prop} then stops at the helper and never unfolds it again.

\begin{lstlisting}[float=t,caption={A helper with its composable lemma (\texttt{Valiant/Gadget.lean}). The attribute \texttt{@[fun\_prop]} registers the second theorem, so that \texttt{fun\_prop} treats \texttt{vv} as a primitive.},label={lst:vv}]
theorem vv_polyTime : ComputableTidilyInPolyTime (fun p : ℕ × Bool => vv p.1 p.2)
    (encPair encUnary encBool) encUnary := by
  unfold vv
  fun_prop

@[fun_prop]
theorem vv_polyTime' {α : Type} {input : α ↪ List Bool} {f : α → ℕ} {g : α → Bool}
    (hf : ComputableTidilyInPolyTime f input encUnary)
    (hg : ComputableTidilyInPolyTime g input encBool) :
    ComputableTidilyInPolyTime (fun a => vv (f a) (g a)) input encUnary :=
  vv_polyTime.comp (hf.prod hg)
\end{lstlisting}

% [rev D1]
Automation stops where the size of intermediate values matters. Proofs by hand were needed for
the binary product of the domain sizes (a fold whose intermediate products are bounded by the
input, through \lstinline{foldlWithInput_of_bounded_growth}), the enumeration of candidate
states (by induction on~$k$), the explicit bounds, and the arithmetic of Valiant's
post-processing.

\subsection{An instance-dependent time bound}
\label{sec:instance-bounds}

% [rev D1]
The bound $c\,(|x|+1)^{en+e}$ of \Cref{sec:boundary} depends on the instance through its number
of attributes~$n$, which is not a function of $|x|$ alone. It is built in three steps.
\begin{enumerate}
\item \emph{A class of bounds.} \lstinline{XPBounded size n T} says that
  $T(a) \le c\,(\mathit{size}(a)+1)^{e\,(n(a)+1)}$ for some constants (\Cref{lst:xp}). It is
  closed under sums, constant factors and fixed powers, and constant factors are absorbed into
  the exponent, since $|x| \ge 1$ (\lstinline{mul_pow_le_pow_add}).
\item \emph{An explicitly bounded enumeration.} All candidate states of $n$ lists are built by a
  left fold, \lstinline{allSections}, whose intermediate states have encodings of length at most
  $4\,(|L|+1)^{n+2}$ (\lstinline{length_state_le}); Complexitylib's fold lemma for
  \lstinline{ComputableTidilyInO} then gives \lstinline{allSections_inO}. Each polynomial-time
  piece enters with a bound $(|x|+1)^{e}$ (\lstinline{inO_pow_of_polyTime}), and
  \lstinline{pipeline_inO} composes a polynomial-time list builder, the enumeration and a
  polynomial-time test into one \lstinline{ComputableTidilyInO} statement with
  \lstinline{XPBounded} time and space.
\item \emph{Explicit constants.} \lstinline{timeAndSpace_of_inO_xp} turns such a statement into
  \lstinline{ComputableTidilyInTimeAndSpace} with the bound $c\,(|x|+1)^{e n + e}$, the form of
  the theorems.
\end{enumerate}

\begin{lstlisting}[float=t,caption={The class of bounds, and the step to explicit constants (\texttt{ExplicitDegree.lean}; statements only). The hypothesis \texttt{h} has a hidden constant factor; the conclusion has none.},label={lst:xp}]
def XPBounded {α : Type} (size n T : α → ℕ) : Prop :=
  ∃ c e, ∀ a, T a ≤ c * (size a + 1) ^ (e * (n a + 1))

theorem timeAndSpace_of_inO_xp {α β : Type} {encIn : α ↪ List Bool} {encOut : β ↪ List Bool}
    {f : α → β} {T S n : α → ℕ} (h : ComputableTidilyInO f encIn encOut T S)
    (hT : XPBounded (fun a => (encIn a).length) n T)
    (hS : XPBounded (fun a => (encIn a).length) n S) :
    ∃ c d, ComputableTidilyInTimeAndSpace f encIn encOut
      (fun a => c * ((encIn a).length + 1) ^ (d * n a + d))
      (fun a => c * ((encIn a).length + 1) ^ (d * n a + d))
\end{lstlisting}

\subsection{Reusable lemmas}
\label{sec:reusable}

% [rev D1]
Some results are independent of decision tables and could serve other developments on
Complexitylib, or be offered upstream; apart from \Cref{def:onecall} they use only
Complexitylib's notions.
\begin{itemize}
\item One-call reductions and their compositions with parsimonious ones
  (\lstinline{OneCallReducible}, \lstinline{ParsimoniousReducible.trans_oneCall},
  \lstinline{OneCallSharpPHard.trans_parsimonious}, \lstinline{SharpPHard.oneCallSharpPHard}).
\item One-call \#P-hardness of \#Monotone-2SAT (\lstinline{sharpMonotone2SAT_oneCallSharpPHard}),
  with the parsimonious encoding of the models of any CNF as the independent sets of size $K$ of
  a graph (\lstinline{modelCount_eq_card}) and the digit lemma for sums of powers
  (\lstinline{digit}).
\item The explicit-bound toolkit: \lstinline{XPBounded} with its closure lemmas,
  \lstinline{inO_pow_of_polyTime}, \lstinline{timeAndSpace_of_inO_xp}, and the bounded
  enumeration of all sections of a list of lists (\lstinline{allSections_inO}).
\item Size facts for encodings, such as \lstinline{length_encList_le_of_inDomains}: a list
  bounded entrywise by another is encoded in no more bits.
\end{itemize}

\subsection{How an AI system was used}
\label{sec:ai}

% [rev D1] [rev m10] Replaces the unnumbered section "Use of generative AI". Springer Nature's
% policy asks that the use of a large language model be documented in the Methods section or a
% suitable alternative part of the manuscript; this subsection is that part.
The development and this paper were written with an AI system, Claude (Anthropic), used through
Claude Code with the model Claude Opus~5.5, between 28~October and 1~November 2026. This subsection records how, from the
repository's history; JAR asks that such use be documented in the methods of the paper.
\begin{itemize}
\item \emph{Who wrote what.} All 80 commits of the repository up to the revision of this paper
  were made by the AI system, in ten sessions, each commit message naming its session. Every
  Lean file of \texttt{LatcComplexity/}, every document of the repository and the text of this
  paper were drafted by it; \texttt{LatcFormal/} is copied from latc-check.
  The author set the questions, scope and venue, directed every session and approved every
  sentence of the final text; the Lean code was written by the AI system.
\item \emph{Briefs.} Each part of the development was done in its own session from a written task description: the initial development (one session, from the Boolean
  model to the conditional monotone case), the binary index with the tractable side, each of the
  three results of the plan \texttt{docs/gaps-plan.md}, the monotone count for arbitrary tables,
  the tightened shape tests, and the shape lemma of \Cref{sec:decision}. Each task description was drafted by an AI assistant at the
  author's direction and approved by the author before its session; the plan and the task
  descriptions of the sessions started that way are in the repository (\texttt{docs/briefs/}),
  and sessions the author steered directly are recorded in their commit messages.
\item \emph{Review of definitions.} Lean checks every proof, so the human review that matters is
  the review of the trusted definitions, which \Cref{app:definitions} collects.
  The definitions were checked against latc-check's model and the README by AI referee reads,
  and the author reviewed their statements. One defect of a definition was found and fixed: the first monotone shape
  tests accepted conditions on attributes that do not exist, so that a clause $[(0,1),(0,1)]$
  over zero attributes passed the test with count~$0$ and the statement ``the all-ones state is
  valid'' failed; the tests now require every attribute index to be declared, and no theorem
  statement changed. It was found during the development.
\item \emph{Iterations on the paper.} The commit messages record four referee reports on an
  early draft, an independent audit and two further referee reads, each followed by a revision,
  and a referee read of the JAR version that led to the present revision. All of these reports
  were produced by AI agents; none was written by a person. The AI system also
  searched for related work and checked bibliographic records against the publishers' records.
\item \emph{Accountability.} The author takes responsibility for the content of the paper and of
  the development. An AI system does not meet the criteria for authorship and is not an author.
  [AUTHOR TO APPROVE]
\end{itemize}

%-----------------------------------------------------------------------------------------------
\section{Conclusion}
\label{sec:conclusion}

% [rev D7]
For decision tables whose rows are complete tuples, completeness under clause constraints is
coNP-complete, finding a missing combination NP-complete, and counting them \#P-complete, with
machine-checked proofs on Turing machines that apply to latc-check's set-based definitions and
hold already for Boolean domains and empty tables. All three are polynomial without
constraints; they are also polynomial with a fixed number of attributes, with a degree linear in
that number. For monotone two-literal constraints, counting the missing combinations stays
\#P-complete under one-call reductions, on tables with arbitrary rows and already with none, by
a new formal proof of a one-call version of Valiant's theorem in the same model. What remains
open in the formalisation is the per-$k$ machine with constants independent of~$k$. The one-call
reductions and the Valiant reduction use only Complexitylib's notions and could be offered
upstream.

%-----------------------------------------------------------------------------------------------
\backmatter

% [rev D4]
\bmhead{Supplementary information}
The Lean development, latc-complexity, described in \Cref{app:artifact}.

% [rev SD]
\section*{Statements and Declarations}
% Headings from the Springer Nature template's Declarations list. Every statement is a
% placeholder for the author, prefilled only from text that already exists in the paper or
% the repository.

\bmhead{Funding}
[AUTHOR TO APPROVE: No funding was received for this work.]

\bmhead{Competing interests}
[AUTHOR TO APPROVE: The author is a co-author of the 2021 paper that introduced the LATC
decision-table completion method, which the author's earlier work~\cite{PaperA} re-examines.
The author has no other competing interests to declare.] (The 2021 paper is~\cite{RaggioRussell2021}.)

\bmhead{Ethics approval and consent to participate}
[AUTHOR TO APPROVE: Not applicable.]

\bmhead{Consent for publication}
[AUTHOR TO APPROVE: Not applicable.]

% [rev D4]
\bmhead{Data availability}
[AUTHOR TO APPROVE: No data were generated or analysed beyond the Lean source code described
under Code availability. The decision-table model it builds on is archived as latc-check~0.1.0
(\url{https://doi.org/10.5281/zenodo.23646699}).]

\bmhead{Materials availability}
[AUTHOR TO APPROVE: Not applicable.]

% [rev D4]
\bmhead{Code availability}
[AUTHOR TO APPROVE: The Lean development latc-complexity, described in \Cref{app:artifact}, will
be released as version~1.0 under the MIT licence at
\url{https://github.com/carlosalbertorussell/latc-complexity} and archived on Zenodo with the
DOI [DOI at release]. The release follows the steps of \texttt{RELEASE.md} in the repository.]

\bmhead{Author contribution}
[AUTHOR TO APPROVE: C.A.R. is the sole author, set the research questions and the task
descriptions, reviewed the definitions and the text, and is responsible for the whole work;
the Lean proofs and the text were drafted with an AI system, as described in \Cref{sec:ai}.]

\bmhead{Use of generative AI}
[AUTHOR TO APPROVE: The use of an AI system is documented in \Cref{sec:ai}.]

%-----------------------------------------------------------------------------------------------
\begin{appendices}

%-----------------------------------------------------------------------------------------------
\section{Main Lean statements}
\label{app:statements}

The statements below are copied from the source up to \lstinline{:=}. Unless noted, they are in
the namespace \lstinline{Latc.PaperD.General}; \lstinline{Inst} is the type of instances of
\Cref{sec:encoding}.

\subparagraph*{Bridge (namespace \lstinline{Latc.PaperD}).}
\begin{lstlisting}
theorem holdsAll_eq_eval (S : Prop5.ConSet ℕ) (x : ℕ → Bool) :
    Prop5.holdsAll S x = (toCNF S).eval x
theorem satisfiable_decode_iff (S : Prop5.ConSet ℕ) :
    Prop5.Satisfiable (Prop5.decode S) ↔ (toCNF S).Satisfiable
\end{lstlisting}

\subparagraph*{Decision problems and counting.}
\begin{lstlisting}
theorem MissingNonempty_npComplete_bin : NPComplete encInstBin MissingNonempty
theorem Complete_coNPComplete_bin : coNPComplete encInstBin Complete
theorem missingNonempty_finset_npComplete_bin :
    NPComplete encInstBin {I : Inst | (missing I.con I.table).Nonempty}
theorem complete_finset_coNPComplete_bin :
    coNPComplete encInstBin {I : Inst | missing I.con I.table = ∅}
theorem missingCount_eq_acceptCount (I : Inst) :
    missingCount I = acceptCount cverifier I (certLen I)
theorem sharpSAT_eq_missingCount (φ : CNF ℕ) : SharpSAT φ = missingCount (fromSAT φ)
theorem missingCount_sharpPComplete_bin : SharpPComplete encInstBin missingCount
\end{lstlisting}

% [rev D5]
\subparagraph*{Boolean domains with empty tables.}
\begin{lstlisting}
theorem fromSAT_mem_boolEmpty (φ : CNF ℕ) : fromSAT φ ∈ BoolEmpty
theorem MissingNonemptyBoolEmpty_npHard_bin : NPHard encInstBin MissingNonemptyBoolEmpty
theorem CompleteBoolEmpty_coNPHard_bin : coNPHard encInstBin CompleteBoolEmpty
theorem missingNonemptyPromise_npHard_bin :
    PromiseProblem.NPHard encInstBin missingNonemptyPromise
theorem missingCountBoolEmpty_sharpPHard_bin : SharpPHard encInstBin missingCountBoolEmpty
\end{lstlisting}

\subparagraph*{One-call reductions and Valiant.}
\lstinline{OneCallReducible} is in \lstinline{Latc.PaperD.General}; the four theorems after it are in
\lstinline{Latc.PaperD.Valiant}.
\begin{lstlisting}
def OneCallReducible (encA : α ↪ List Bool) (encB : β ↪ List Bool) (f : α → ℕ) (g : β → ℕ) :
    Prop :=
  ∃ (r : α → β) (post : α → ℕ → ℕ), ComputableInPolyTime r encA encB ∧
    ComputableInPolyTime (fun p : α × ℕ => post p.1 p.2) (encPair encA encNat) encNat ∧
      ∀ a, f a = post a (g (r a))
theorem sharpMonotone2SAT_oneCallSharpPHard :
    OneCallSharpPHard encCNF SharpMonotone2SAT
theorem countValidMonotone2_sharpPComplete_oneCall :
    OneCallSharpPComplete encInstBin countValidMonotone2
theorem countMissingMonotone2_eq_card_missing (I : General.Inst) (h : pos2AnyB I = true) :
    countMissingMonotone2 I = (Latc.missing I.con I.table).card
theorem countMissingMonotone2_sharpPComplete_oneCall :
    OneCallSharpPComplete encInstBin countMissingMonotone2
\end{lstlisting}

\subparagraph*{Tractability.}
\begin{lstlisting}
theorem missingCount_eq_of_noConstraints (I : Inst) (h : I.2.1 = []) :
    missingCount I = domProd I - distinctRows I
theorem missingCountFixedAttrs_polyTime (k : ℕ) :
    ComputableInPolyTime (fun I : Inst => if I.n ≤ k then missingCount I else 0)
      encInstBin encNat
theorem missingNonempty_explicitDegree :
    ∃ c d, ComputableTidilyInTimeAndSpace
      (fun I : Inst => decide (missing I.con I.table).Nonempty) encInstBin encBool
      (fun I => c * ((encInstBin I).length + 1) ^ (d * I.n + d))
      (fun I => c * ((encInstBin I).length + 1) ^ (d * I.n + d))
theorem missingNonemptyFixedAttrs_uniform :
    ∃ c d₀ d₁, ∀ k, ComputableTidilyInTimeAndSpace
      (fun I : {I : Inst // I.n ≤ k} => decide (missing I.1.con I.1.table).Nonempty)
      ((Function.Embedding.subtype _).trans encInstBin) encBool
      (fun I => c * ((encInstBin I.1).length + 1) ^ (d₀ * k + d₁))
      (fun I => c * ((encInstBin I.1).length + 1) ^ (d₀ * k + d₁))
\end{lstlisting}

% [rev m6]
\subparagraph*{The unary forms.}
Each polynomial-time result of \Cref{sec:boundary} also holds for \lstinline{encInst}:
\lstinline{completeNoConstraints_mem_P_unary}, \lstinline{missingNonemptyNoConstraints_mem_P_unary},
\lstinline{missingCountNoConstraints_polyTime_unary}, \lstinline{completeFixedAttrs_mem_P_unary},
\lstinline{missingNonemptyFixedAttrs_mem_P_unary} and \lstinline{missingCountFixedAttrs_polyTime_unary},
the first five through \lstinline{mem_P_of_mem_P_bin} or the re-encoding of \Cref{sec:encoding}.
The hardness results have their unary forms without the suffix \lstinline{_bin}
(\lstinline{MissingNonempty_npComplete}, \lstinline{Complete_coNPComplete},
\lstinline{missingNonempty_finset_npComplete}, \lstinline{complete_finset_coNPComplete},
\lstinline{missingCount_sharpPComplete}), and the monotone ones with the suffix
\lstinline{_unary} (\lstinline{countValidMonotone2_sharpPComplete_oneCall_unary},
\lstinline{countMissingMonotone2_sharpPComplete_oneCall_unary}).

%-----------------------------------------------------------------------------------------------
% [rev D2] New appendix.
\section{Trusted definitions}
\label{app:definitions}

The definitions below are the ones a reader must trust beyond Complexitylib, CSLib and Mathlib
(\Cref{sec:formal}). They are copied verbatim from the source, with their documentation
comments, file by file. Namespaces: \lstinline{Latc} for latc-check's model,
\lstinline{Latc.PaperD.General} for the rest except \lstinline{pos2AnyB} and
\lstinline{countMissingMonotone2}, which are in \lstinline{Latc.PaperD.Valiant}.

\subparagraph*{latc-check's model (\texttt{LatcFormal/Basic.lean}).}
Here \lstinline{ι} is a finite type of attributes, \lstinline{D i} a finite domain for each, and
\lstinline{C} a decidable predicate on tuples.
\begin{lstlisting}
/-- The valid space 𝒱 = {x ∈ 𝒟 : c(x) for all c ∈ C} (Section 3.1). -/
def validSpace : Finset (∀ i, D i) := univ.filter C

/-- The missing set M(T₀) = 𝒱 ∖ cond(T₀) (Section 3.2). -/
def missing (T : Finset (∀ i, D i)) : Finset (∀ i, D i) := validSpace C \ T
\end{lstlisting}

\subparagraph*{Instances and the decision problems (\texttt{General.lean}).}
\begin{lstlisting}
/-- A condition Aᵢ = v, as the pair `(i, v)`. -/
abbrev Cond := ℕ × ℕ

/-- An instance: domain sizes, clause constraints, and the condition tuples of T₀. -/
abbrev Inst := List ℕ × (List (List Cond) × List (List ℕ))

/-- Instances: binary domain sizes, values and rows; unary attribute indices. -/
abbrev encInst : Inst ↪ List Bool :=
  encPair (encList encNat) (encPair (encList (encList (encPair encUnary encNat)))
    (encList (encList encNat)))

/-- `s` is a state of the declared domains `d`: one value per attribute, inside its domain. -/
def InDomains (d s : List ℕ) : Prop := List.Forall₂ (· < ·) s d

/-- The state `s` satisfies the condition Aᵢ = v. -/
def CondHolds (s : List ℕ) (a : Cond) : Prop := s[a.1]? = some a.2

/-- The state `s` satisfies every clause constraint of `S`. -/
def Valid (S : List (List Cond)) (s : List ℕ) : Prop := ∀ c ∈ S, ∃ a ∈ c, CondHolds s a

/-- M(T₀) is non-empty: some valid state of the declared domains is not a condition tuple
of T₀. -/
def MissingNonempty : Set Inst :=
  {I | ∃ s, InDomains I.1 s ∧ Valid I.2.1 s ∧ s ∉ I.2.2}

/-- T₀ is complete: M(T₀) = ∅. -/
def Complete : Set Inst := MissingNonemptyᶜ
\end{lstlisting}

\subparagraph*{The binary encoding (\texttt{Binary.lean}).}
\begin{lstlisting}
/-- Instances with every number in binary, the attribute index of a condition included. -/
abbrev encInstBin : Inst ↪ List Bool :=
  encPair (encList encNat) (encPair (encList (encList (encPair encNat encNat)))
    (encList (encList encNat)))
\end{lstlisting}

\subparagraph*{The table an instance denotes (\texttt{Adequacy.lean}).}
\begin{lstlisting}
/-- The number of attributes. -/
abbrev Inst.n (I : Inst) : ℕ := I.1.length

/-- The declared domain of attribute `i`: the values `0, …, dᵢ - 1`. -/
abbrev Inst.Dom (I : Inst) (i : Fin I.n) : Type := Fin (I.1.get i)

/-- The list of values of a condition tuple. -/
def Inst.toList (I : Inst) (x : ∀ i, I.Dom i) : List ℕ := List.ofFn fun i => (x i).val

/-- The constraint predicate of the instance, on condition tuples. -/
def Inst.con (I : Inst) (x : ∀ i, I.Dom i) : Prop := Valid I.2.1 (I.toList x)

/-- The condition set of T₀: the tuples whose list of values is a row of `T`. -/
def Inst.table (I : Inst) : Finset (∀ i, I.Dom i) :=
  Finset.univ.filter fun x => I.toList x ∈ I.2.2
\end{lstlisting}

\subparagraph*{Counting (\texttt{CountingGeneral.lean}).}
\begin{lstlisting}
def missingCount (I : Inst) : ℕ := (missing I.con I.table).card
\end{lstlisting}

\subparagraph*{One-call reductions and the monotone problems (\texttt{Monotone.lean}).}
\begin{lstlisting}
/-- `f` reduces to `g` with one oracle call: `f a = post a (g (r a))` for polynomial-time `r` and
`post`, the oracle answer read in binary. -/
def OneCallReducible (encA : α ↪ List Bool) (encB : β ↪ List Bool) (f : α → ℕ) (g : β → ℕ) :
    Prop :=
  ∃ (r : α → β) (post : α → ℕ → ℕ), ComputableInPolyTime r encA encB ∧
    ComputableInPolyTime (fun p : α × ℕ => post p.1 p.2) (encPair encA encNat) encNat ∧
      ∀ a, f a = post a (g (r a))

/-- Every function in `#P`, with any encoding of its instances, reduces to `g` with one call. -/
def OneCallSharpPHard (encB : β ↪ List Bool) (g : β → ℕ) : Prop :=
  ∀ ⦃α : Type⦄ (enc : α ↪ List Bool) (f : α → ℕ), f ∈ SharpP enc → OneCallReducible enc encB f g

/-- `g` is in `#P` and one-call `#P`-hard. -/
def OneCallSharpPComplete (encB : β ↪ List Bool) (g : β → ℕ) : Prop :=
  g ∈ SharpP encB ∧ OneCallSharpPHard encB g

/-- Every clause has two literals, both positive: `xᵢ ∨ xⱼ`. -/
def monotone2B (φ : CNF ℕ) : Bool := φ.all fun c => c.length == 2 && c.all (·.2)

/-- `#Monotone-2SAT`: the model count of a monotone two-literal formula; zero on other syntax. -/
def SharpMonotone2SAT (φ : CNF ℕ) : ℕ := if monotone2B φ then SharpSAT φ else 0

/-- Domains {0, 1}, clauses of two conditions with value 1 on declared attributes (index below
the number of attributes), and T₀ = ∅. -/
def pos2B (I : Inst) : Bool :=
  I.1.all (· == 2) &&
      I.2.1.all (fun c =>
        c.length == 2 && c.all (fun a => decide (a.1 < I.1.length) && a.2 == 1)) &&
    I.2.2.isEmpty

/-- |𝒱| for clause constraints of two conditions Aᵢ = 1, Aⱼ = 1 over domains {0, 1}, in
latc-check's model (`missingCount` with T₀ = ∅); zero on other syntax. -/
noncomputable def countValidMonotone2 (I : Inst) : ℕ := if pos2B I then missingCount I else 0
\end{lstlisting}

\subparagraph*{The monotone count for any table (\texttt{Valiant/AnyTable.lean}).}
\begin{lstlisting}
/-- Domains {0, 1} and clauses of two conditions with value 1 on declared attributes (index below
the number of attributes); any table T₀. -/
def pos2AnyB (I : General.Inst) : Bool :=
  I.1.all (· == 2) &&
    I.2.1.all (fun c =>
      c.length == 2 && c.all (fun a => decide (a.1 < I.1.length) && a.2 == 1))

/-- |M(T₀)| for clause constraints of two conditions Aᵢ = 1, Aⱼ = 1 over domains {0, 1}, any
table T₀, in latc-check's model; zero on other syntax. -/
noncomputable def countMissingMonotone2 (I : General.Inst) : ℕ :=
  if pos2AnyB I then missingCount I else 0
\end{lstlisting}

\subparagraph*{Boolean domains with empty tables (\texttt{Restricted.lean}).}
\begin{lstlisting}
/-- Boolean domains and an empty table: every declared domain size is 2 and T₀ has no rows. -/
def BoolEmpty : Set Inst := {I | (∀ d ∈ I.1, d = 2) ∧ I.2.2 = []}

/-- M(T₀) ≠ ∅, on Boolean domains with T₀ = ∅. -/
def MissingNonemptyBoolEmpty : Set Inst := MissingNonempty ∩ BoolEmpty

/-- M(T₀) = ∅, on Boolean domains with T₀ = ∅. -/
def CompleteBoolEmpty : Set Inst := Complete ∩ BoolEmpty

/-- The promise problem: on Boolean domains with T₀ = ∅, tell M(T₀) ≠ ∅ from M(T₀) = ∅. -/
def missingNonemptyPromise : PromiseProblem Inst :=
  (PromiseProblem.ofSet MissingNonempty).restrict BoolEmpty

/-- |M(T₀)| on Boolean domains with T₀ = ∅, and `0` on other instances. -/
def missingCountBoolEmpty (I : Inst) : ℕ := if I ∈ BoolEmpty then missingCount I else 0
\end{lstlisting}

%-----------------------------------------------------------------------------------------------
\section{Artifact and reproducibility}
\label{app:artifact}

\subparagraph*{Contents.}
% [rev D1] [rev m13]
The development is the directory \texttt{LatcComplexity/} of the repository latc-complexity:
4\,903 lines of Lean in 17 files, 4\,520 without the two axiom-audit files
(\Cref{tab:modules}), on top of latc-check's model in \texttt{LatcFormal/}. The eleven files of
\texttt{LatcFormal/} and its root file are byte-identical to the folder \texttt{formal/} of the
archived release latc-check~0.1.0~\cite{latc_check_0_1_0}. The README lists every main theorem
with its statement, and \texttt{docs/paper-facts.md} records the fully qualified name, file and
line of every name this paper cites.

% [rev m8]
\subparagraph*{The Boolean model.}
The repository also contains the first form of the development, a Boolean model
(\texttt{Bridge.lean}, \texttt{Counting.lean}): its attributes are the variables of a
Complexitylib CNF formula used as the constraints, and its rows are bit vectors. NP-, coNP- and
\#P-completeness are proved for it too (\lstinline{Latc.PaperD.MissingNonempty_npComplete},
\lstinline{Latc.PaperD.Complete_coNPComplete}, \lstinline{Latc.PaperD.countMissing_sharpPComplete}),
the last with a parsimonious reduction \emph{to} \#SAT (\lstinline{countMissing_eq_sharpSAT}).
It has no adequacy theorem, and the results of this paper do not use it.

% [rev D4] [rev m9]
\subparagraph*{Versions.}
Lean \texttt{v4.36.0-rc1}; Complexitylib at commit \texttt{b7acfc2} (28 October 2026), with the
CSLib and Mathlib it pins. The earlier, pre-CSLib version of Complexitylib, which
Russell~\cite{PaperA} cites, is a separate repository with a different API.

% [rev D4]
\subparagraph*{Build and audit.}
The build reported here is a cold build of commit \texttt{0ea3170}, whose Lean files are
those of this paper, made on 1 November 2026 on a Linux x86-64 machine with four cores, in a
fresh clone with no \texttt{\textasciitilde/.elan}, \texttt{.lake} or Mathlib cache, with Lean
\texttt{4.36.0-rc1}, Lake 5.0.0 and elan 4.2.4. \texttt{lake exe cache get} downloaded the
toolchain and all 9\,047 precompiled Mathlib files in 103~s. \texttt{lake build} compiled 239
modules of CSLib and Complexitylib and the 30 modules of the repository in 217~s and ended with
\texttt{Build completed successfully (18651 jobs)}, exit code~0, no warnings and no errors; the
checkout, toolchain and caches took 11.5~GB of disk. The axiom audit, \texttt{lake env lean} on
the three audit files (\texttt{LatcComplexity/Axioms.lean}, \texttt{LatcComplexity/Valiant/Axioms.lean}
and \texttt{LatcFormal/Axioms.lean}), took 13~s and printed 397 answers: 344 use all three
standard axioms, 51 use one or two of them, 2 use none, and none uses \lstinline{sorryAx}. A
whole-word search of the \texttt{.lean} files finds \lstinline{sorry}, \texttt{admit},
\texttt{native\_decide} and \texttt{axiom} only in the \texttt{\#print axioms} commands.

% [rev D4]
\subparagraph*{Release.}
The repository is private during review. The file \texttt{RELEASE.md} lists the steps of the
release: tagging version~1.0, a cold build and the full audit at the tag, making the repository
public, and archiving the tag on Zenodo, whose DOI will be cited here [DOI at release]. The
commit-by-commit history of earlier builds is in the README.

\end{appendices}

%-----------------------------------------------------------------------------------------------
\bibliography{../../docs/paper-refs,../extra}

\end{document}
